\documentclass[autodetect-engine,11pt,a4paper]{bxjsarticle}
\usepackage{amsmath,amssymb}
\usepackage{geometry}
\geometry{margin=20mm}

\title{\Large\textbf{高エネルギー物理学 レポート課題}}
\author{BHB26052 河村太星}
\date{\today}

\begin{document}

\maketitle

\section*{第1問：弱い相互作用のファインマン図}
以下の弱い相互作用による過程の最低次(クォークレベル)ファインマン図を示す。

\subsection*{(1) $\nu_e + n \rightarrow e^- + p$}
$W^+$ボソンの交換による電荷カレント(CC)過程である。中性子$n(udd)$内の$d$クォークが$W^-$ボソンを放出して$u$クォークへ変化し(陽子$p(uud)$となる)、$W^-$が$\nu_e$に吸収されて$e^-$になる(あるいは$W^+$の交換として記述される)。

\vspace{5cm} % 手書き・図貼り付け用のスペース

\subsection*{(2) $\pi^+ \rightarrow \mu^+ + \nu_\mu$}
荷電パイ中間子$\pi^+(u\bar{d})$の純レプトン崩壊である。$u$クォークと$\bar{d}$反クォークが対消滅して$W^+$ボソンを生成し、その$W^+$が$\mu^+$と$\nu_\mu$に崩壊する。

\vspace{5cm} % 手書き・図貼り付け用のスペース

\subsection*{(3) $K^0 \rightarrow \pi^+ + \pi^-$}
ストレンジネスの変化を伴う弱崩壊($\Delta S = 1$)である。$K^0(d\bar{s})$中の$\bar{s}$反クォークが$W^+$を放出して$\bar{u}$反クォークに変化し、$W^+$が$u\bar{d}$ペアに崩壊する。残った$d$クォークは観測線としてそのまま残り、それぞれ$\pi^+(u\bar{d})$と$\pi^-(d\bar{u})$を形成する。

\vspace{5cm} % 手書き・図貼り付け用のスペース

\newpage

\section*{第2問：強い相互作用のクォークフロー図}
以下の強い相互作用による反応または崩壊のクォークフロー図を示す。

\subsection*{(1) $p + \pi^- \rightarrow n + \pi^0$}
$p(uud)$と$\pi^-(d\bar{u})$の反応により、$n(udd)$と$\pi^0\left(\frac{u\bar{u}-d\bar{d}}{\sqrt{2}}\right)$が生成される。$p$の$u$と$\pi^-$の$\bar{u}$が消滅し、グルーオンを介して$d\bar{d}$ペア(あるいは$u\bar{u}$ペア)が再結合する。

\vspace{5cm} % 手書き・図貼り付け用のスペース

\subsection*{(2) $p + K^- \rightarrow K^+ + \Xi^-$}
$p(uud)$と$K^-(s\bar{u})$ の反応である。強い相互作用によりグルーオンから$s\bar{s}$ペアが生成され、$K^-(s\bar{u})$ の $s$ クォークとともに $\Xi^-(dss)$ および $K^+(u\bar{s})$ が形成される。

\vspace{5cm} % 手書き・図貼り付け用のスペース

\subsection*{(3) $\Delta^{++} \rightarrow p + \pi^+$}
$\Delta^{++}(uuu)$ の強い崩壊である。グルーオン放射を介して $d\bar{d}$ ペアが生成され、$uuu$ のうち $uu$ と生成された $d$ が結合して $p(uud)$ に、$u$ と $\bar{d}$ が結合して $\pi^+(u\bar{d})$ になる。

\vspace{5cm} % 手書き・図貼り付け用のスペース


\section*{第3問：アイソスピン保存則と断面積の比}

各粒子のアイソスピン状態は以下の通りである。
\begin{align*}
    \text{パイ中間子 } (I=1): \quad & |\pi^+\rangle = |1, 1\rangle, \quad |\pi^0\rangle = |1, 0\rangle, \quad |\pi^-\rangle = |1, -1\rangle \\
    \text{核子 } (I=1/2): \quad & |p\rangle = |1/2, 1/2\rangle, \quad |n\rangle = |1/2, -1/2\rangle \\
    \text{K中間子 } (I=1/2): \quad & |K^+\rangle = |1/2, 1/2\rangle, \quad |K^0\rangle = |1/2, -1/2\rangle \\
    \text{シグマバリオ } (I=1): \quad & |\Sigma^+\rangle = |1, 1\rangle, \quad |\Sigma^0\rangle = |1, 0\rangle, \quad |\Sigma^-\rangle = |1, -1\rangle
\end{align*}

Clebsch-Gordan係数を用いて、始状態および終状態を全アイソスピン状態$|I, I_3\rangle$に分解する。

\subsection*{始状態の分解}
\begin{align*}
    |\pi^- p\rangle &= |1, -1\rangle \otimes |1/2, 1/2\rangle = \sqrt{\frac{1}{3}}|3/2, -1/2\rangle - \sqrt{\frac{2}{3}}|1/2, -1/2\rangle \\
    |\pi^+ p\rangle &= |1, 1\rangle \otimes |1/2, 1/2\rangle = |3/2, 3/2\rangle
\end{align*}

\subsection*{終状態の分解}
\begin{align*}
    |K^0 \Sigma^0\rangle &= |1/2, -1/2\rangle \otimes |1, 0\rangle = \sqrt{\frac{2}{3}}|3/2, -1/2\rangle + \sqrt{\frac{1}{3}}|1/2, -1/2\rangle \\
    |K^+ \Sigma^-\rangle &= |1/2, 1/2\rangle \otimes |1, -1\rangle = \sqrt{\frac{1}{3}}|3/2, -1/2\rangle - \sqrt{\frac{2}{3}}|1/2, -1/2\rangle \\
    |K^+ \Sigma^+\rangle &= |1/2, 1/2\rangle \otimes |1, 1\rangle = |3/2, 3/2\rangle
\end{align*}

\subsection*{遷移振幅の計算}
アイソスピン保存により、各反応の行列要素$M^{(i)}$は$I=3/2, 1/2$の振幅$M_3, M_1$を用いて表される。
\begin{align*}
    M^{(1)} &= \langle K^0 \Sigma^0 | \mathcal{M} | \pi^- p \rangle = \sqrt{\frac{2}{3}}\sqrt{\frac{1}{3}}M_3 - \sqrt{\frac{1}{3}}\sqrt{\frac{2}{3}}M_1 = \frac{\sqrt{2}}{3}(M_3 - M_1) \\
    M^{(2)} &= \langle K^+ \Sigma^- | \mathcal{M} | \pi^- p \rangle = \sqrt{\frac{1}{3}}\sqrt{\frac{1}{3}}M_3 + \left(-\sqrt{\frac{2}{3}}\right)\left(-\sqrt{\frac{2}{3}}\right)M_1 = \frac{1}{3}M_3 + \frac{2}{3}M_1 \\
    M^{(3)} &= \langle K^+ \Sigma^+ | \mathcal{M} | \pi^+ p \rangle = M_3
\end{align*}

\subsection*{断面積の比}
断面積は$\sigma \propto |M|^2$より、求める比は以下のようになる。
\[
\sigma(\pi^- p \rightarrow K^0 \Sigma^0) : \sigma(\pi^- p \rightarrow K^+ \Sigma^-) : \sigma(\pi^+ p \rightarrow K^+ \Sigma^+) = \frac{2}{9}|M_3 - M_1|^2 : \frac{1}{9}|M_3 + 2M_1|^2 : |M_3|^2
\]


\section*{第4問：相対論的運動学とローレンツ変換}

\subsection*{1. 重心系$S'$における4元運動量}
$\pi^0$中間子の静止系(重心系$S'$)において、$\pi^0$の質量を$m$とする。崩壊により生じる2つの光子$\gamma_1, \gamma_2$は等しいエネルギー$E' = m/2$を持ち、進行方向($x'$軸)に対して垂直($y'$軸方向)に放出されたため、4元運動量は次のように表される。
\[
p_{\gamma 1}^{\prime \mu} = \left( \frac{m}{2}, 0, \frac{m}{2}, 0 \right), \quad p_{\gamma 2}^{\prime \mu} = \left( \frac{m}{2}, 0, -\frac{m}{2}, 0 \right)
\]

\subsection*{2. ローレンツ変換による実験室系$S$への変換}
$\pi^0$が実験室系$S$で$x$軸方向に運動量$p = 1.0\text{ GeV/c}$で運動しているとする。ローレンツ因子は
\[
E = \sqrt{p^2 + m^2}, \quad \gamma = \frac{E}{m}, \quad \beta = \frac{p}{E}
\]
である。$x$軸方向のローレンツ変換式
\[
E_i = \gamma (E'_i + \beta p'_{xi}), \quad p_{xi} = \gamma (p'_{xi} + \beta E'_i), \quad p_{yi} = p'_{yi}
\]
を適用すると、光子1の実験室系での運動量は
\[
E_1 = \gamma \frac{m}{2} = \frac{E}{2}, \quad p_{x1} = \gamma \beta \frac{m}{2} = \frac{p}{2}, \quad p_{y1} = \frac{m}{2}
\]
となる。

\subsection*{3. 2光子の放出角度$\theta$}
光子1が$x$軸(元の$\pi^0$の進行方向)となす角を$\theta/2$とすると、
\[
\tan\left(\frac{\theta}{2}\right) = \frac{p_{y1}}{p_{x1}} = \frac{m/2}{p/2} = \frac{m}{p}
\]
したがって、2つの光子の運動方向がなす角度$\theta$は
\[
\theta = 2 \arctan\left(\frac{m}{p}\right)
\]
となる。

$\pi^0$の質量$m \approx 0.135\text{ GeV/c}^2$および$p = 1.0\text{ GeV/c}$を代入すると、
\[
\frac{m}{p} = 0.135 \implies \frac{\theta}{2} = \arctan(0.135) \approx 0.1342\text{ rad} \quad (7.69^\circ)
\]
\[
\theta \approx 0.268\text{ rad} \quad (15.4^\circ)
\]
したがって、実験室系で見た2つの光子の運動方向がなす角度は約$15.4^\circ$($0.268\text{ rad}$)である。

\end{document}